\documentclass[unnumsec,webpdf,contemporary,large]{oup-authoring-template}
%% ---------------------------------------------------------------------------
%% Glaukopis: Cyber Threat Intelligence Language Model Training:
%% Optimizing Domain Specific LLM Performance
%%
%% Worth, Cardei, Alam. Formatted for the OUP contemporary
%% author-year template. Compile with pdfLaTeX.
%%
%% NOTE: This template already loads amsmath, amssymb, hyperref, graphicx,
%% url, natbib, etc. Do NOT re-load them.
%% ---------------------------------------------------------------------------

\graphicspath{{Fig/}}

%% Theorem-style declarations (OUP-standard; kept for completeness even if unused)
\theoremstyle{thmstyleone}
\newtheorem{theorem}{Theorem}
\newtheorem{proposition}[theorem]{Proposition}
\theoremstyle{thmstyletwo}
\newtheorem{example}{Example}
\newtheorem{remark}{Remark}
\theoremstyle{thmstylethree}
\newtheorem{definition}{Definition}

\begin{document}

\journaltitle{Journal of Cybersecurity}
\DOI{DOI added during production}
\copyrightyear{2026}
\pubyear{2026}
\vol{XX}
\issue{x}
\access{Published: Date added during production}
\appnotes{Paper}

\firstpage{1}

\title[Glaukopis: CTI LLM Training]{Glaukopis: Cyber Threat Intelligence Language Model Training: Optimizing Domain Specific LLM Performance}

\author[1,$\ast$]{Peter J. Worth Jr.}
\author[1]{Ionut Cardei}
\author[1]{Md Tanvirul Alam}

\address[1]{\orgdiv{Department of Electrical Engineering and Computer Science}, \orgname{Florida Atlantic University}, \orgaddress{\street{777 Glades Road}, \postcode{33431}, \state{Boca Raton, FL}, \country{United States}}}

\corresp[$\ast$]{Corresponding author. \href{mailto:pworth@athenasecuritygroup.com}{pworth@athenasecuritygroup.com}}

\received{Date}{0}{Year}
\revised{Date}{0}{Year}
\accepted{Date}{0}{Year}

\abstract{Large language models (LLMs) are increasingly applied to cybersecurity tasks ranging from vulnerability analysis to cyber threat intelligence (CTI) reporting. However, generic foundation models often lack the domain knowledge, calibration, and reasoning needed for high-stakes CTI workflows, while frontier models raise cost, latency, and privacy concerns. Recent work from IBM on CyberPal.AI and its follow-on line ``Toward Cybersecurity-Expert Small Language Models'' demonstrates that expert-curated instruction fine-tuning (IFT) can transform small language models (SLMs) into capable cybersecurity assistants when paired with high-quality, domain-specific instruction datasets such as SecKnowledge and SecKnowledge~2.0 \cite{levi2024cyberpal,levi2025toward}.
In parallel, the CTI community has developed rich knowledge bases (MITRE ATT\&CK, CVE, CAPEC, CWE, ENGAGE) and domain-specific benchmarks such as CTIBench and, most recently, AthenaBench---a dynamic, multi-task CTI benchmark that integrates live data sources and emphasizes applied security reasoning \cite{strom2018mitreattck,mitreCVE,mitreCWE,mitreCAPEC,alam2024ctibench,alam2025athenabench}.
This paper surveys the research foundations underlying an ongoing effort at Florida Atlantic University and Athena Labs (Athena Security Group) to build instruction-tuned cybersecurity SLMs using a template-based IFT data generation system grounded in a CTI knowledge graph. We (1) review instruction fine-tuning and SLMs with emphasis on CyberPal and SecKnowledge; (2) describe the role of CTI knowledge bases and knowledge graphs in structuring cyber reasoning tasks; (3) situate Athena's template-based IFT generator within this landscape; and (4) analyze the state of cybersecurity LLM benchmarking, with particular focus on AthenaBench as an emerging standard for CTI evaluation. We conclude by outlining research directions for aligning template-generated IFT corpora, cybersecurity SLMs, and dynamic CTI benchmarks.}

\keywords{instruction fine-tuning; small language models; cyber threat intelligence; CTI knowledge graphs; CyberPal; AthenaBench; reinforcement learning with verifiable rewards}

\keywords[Abbreviations]{ATT\&CK, CAPEC, CoT, CTI, CVE, CWE, IFT, KEV, KG, LLM, NVD, RLHF, SFT, SLM, SOC, TTP}

\maketitle

\section{Introduction}

LLMs have rapidly become integral to many security workflows: summarizing threat reports, mapping indicators to MITRE ATT\&CK, triaging vulnerabilities, and assisting with detection engineering. Surveys of LLMs for security show use cases spanning vulnerability detection, malware analysis, CTI, phishing detection, and SOC automation, but also highlight persistent issues: hallucinations, shallow pattern matching rather than causal reasoning, and limited ability to operate on up-to-date threat information \cite{levi2024cyberpal,liu2024cyberbench}.

Two trends have emerged in response:

\begin{itemize}
  \item \textbf{Domain-specific models and datasets.} Models such as SecureBERT and CySecBERT pretrain or adapt transformers on cybersecurity corpora, improving representation quality for CTI and related tasks. Instruction-tuned cybersecurity assistants such as Lily-Cybersecurity-7B and IBM's CyberPal models further specialize LLMs for interactive security workflows using cyber-focused instruction--response pairs \cite{levi2024cyberpal}.
  \item \textbf{Domain-specific benchmarks.} Benchmarks such as CyberMetric, CyberBench, SecEval, CTIBench, and SEvenLLM-Bench evaluate security knowledge, reasoning, and incident analysis in LLMs \cite{alam2024ctibench,ji2024sevenllm,liu2024cyberbench}.
\end{itemize}

However, many existing benchmarks are static and quickly fall out of sync with the evolving threat landscape. AthenaBench addresses this by introducing a \emph{dynamic CTI benchmark} that continuously draws from live sources (MITRE ATT\&CK, NVD, current APT reports) to construct tasks targeting CTI knowledge, vulnerability severity, root cause mapping, threat actor attribution, mitigation strategies, and attack technique extraction \cite{alam2025athenabench}.

In parallel, the IFT design document by Cardei describes an \emph{IFT Dataset Generation System} that compiles high-level templates into (instruction, question, answer) triples grounded in an ASG Neo4j CTI graph built from MITRE and other CTI sources \cite{cardei2025ift}. This system is intended to produce large, structured, and incrementally refreshable cybersecurity IFT corpora for training SLMs specialized to CTI workflows.

This paper positions that work in its research context:
\begin{itemize}
  \item Section~\ref{sec:ift-slm} reviews instruction fine-tuning and small language models, emphasizing IBM's CyberPal line.
  \item Section~\ref{sec:cti-kg} discusses CTI, CTI knowledge bases, and the role of knowledge graphs and templates in generating IFT data.
  \item Section~\ref{sec:benchmarks} surveys cybersecurity LLM benchmarks, highlighting AthenaBench and its relationship to CTIBench and other efforts.
  \item Section~\ref{sec:athena-ift} connects these strands to the Athena IFT generator design and outlines research opportunities.
\end{itemize}


\section{Instruction Fine-Tuning: History and Context}
\label{sec:ift-history}

Instruction fine-tuning (IFT) has emerged as a foundational paradigm for aligning large language models (LLMs) with human intent and task-oriented usage. Rather than treating language modeling as an end in itself, IFT reframes downstream interaction as the problem of mapping natural-language instructions to correct, task-specific outputs. Although early work on task fine-tuning existed prior to large-scale transformers, the modern formulation of instruction fine-tuning coalesced through three seminal research efforts between 2021 and 2022. Together, these works established the empirical phenomenon, the alignment pipeline, and the scalability properties that underpin nearly all contemporary instruction-tuned models, including domain-specialized systems such as CyberPal.AI and the Athena CTI instruction framework.

\subsection{FLAN: Instruction Tuning as a Generalization Paradigm}

The first widely recognized formulation of instruction fine-tuning was introduced by Wei et al.\ in \emph{Finetuned Language Models Are Zero-Shot Learners} (FLAN) \cite{wei2022flan}. This work demonstrated that large pretrained language models, when fine-tuned on a sufficiently diverse collection of tasks expressed uniformly as natural-language instructions, exhibit strong zero-shot generalization to unseen tasks.

FLAN unified dozens of existing NLP datasets by converting each task into multiple instruction templates (e.g., question answering, sentiment analysis, textual entailment, translation). The authors showed that instruction formatting itself---explicitly describing the task in natural language---was a critical factor in enabling generalization, often rivaling or surpassing few-shot prompting on much larger models.

The central insight of FLAN is that \emph{instruction-following is a learnable capability} rather than an emergent property of scale alone. By training on heterogeneous instruction--response pairs, the model internalizes a latent task distribution that allows it to interpret novel instructions without task-specific fine-tuning. This result established instruction fine-tuning as a distinct and powerful paradigm, forming the conceptual foundation for subsequent work in alignment and domain adaptation.

\subsection{InstructGPT: Instruction Following and Alignment with Human Feedback}

While FLAN established the effectiveness of instruction fine-tuning for generalization, it did not directly address issues of safety, helpfulness, or alignment with human preferences. These concerns were tackled by Ouyang et al.\ in \emph{Training Language Models to Follow Instructions with Human Feedback} (InstructGPT) \cite{ouyang2022instructgpt}.

InstructGPT introduced a three-stage training pipeline that remains the basis for many production LLMs:
\begin{enumerate}
  \item \textbf{Supervised instruction fine-tuning (SFT)} on human-written instruction--response pairs;
  \item \textbf{Preference modeling} using human rankings of multiple model outputs;
  \item \textbf{Reinforcement learning from human feedback (RLHF)} to optimize model behavior toward preferred responses.
\end{enumerate}

A key result of this work was that relatively small models, when instruction-tuned and aligned via RLHF, were preferred by human evaluators over much larger pretrained models. This demonstrated that alignment and instruction-following quality can outweigh raw model scale.

From a historical perspective, InstructGPT complements FLAN by transforming instruction fine-tuning from a purely performance-driven technique into a broader alignment framework. Together, FLAN and InstructGPT established instruction fine-tuning as both a generalization mechanism and a control interface between humans and large language models.

\subsection{FLAN-T5: Scaling Instruction Fine-Tuning}

The third pillar of instruction fine-tuning was established by Chung et al.\ in \emph{Scaling Instruction-Finetuned Language Models} (FLAN-T5) \cite{chung2022flant5}. This work systematically explored how instruction fine-tuning behaves across model sizes, task mixtures, and dataset scales.

FLAN-T5 expanded the instruction corpus to thousands of tasks and demonstrated that instruction fine-tuning consistently improves performance across a wide range of downstream benchmarks, including reasoning and knowledge-intensive tasks. Importantly, the authors showed that instruction fine-tuning benefits models of all sizes, from small to very large, and that gains persist even as base models scale.

This result established instruction fine-tuning as a \emph{scalable and robust training paradigm}, rather than a technique limited to a narrow regime. It also reinforced the idea that the diversity and structure of instruction data are as important as parameter count---a conclusion that directly motivates modern domain-specialized instruction datasets.

\subsection{Implications for Domain-Specialized Instruction Tuning}

Together, FLAN, InstructGPT, and FLAN-T5 define the historical and conceptual foundations of instruction fine-tuning:
\begin{itemize}
  \item FLAN establishes that instruction-following can be learned and generalized;
  \item InstructGPT integrates instruction fine-tuning with human alignment objectives;
  \item FLAN-T5 demonstrates that instruction fine-tuning scales across model sizes and task mixtures.
\end{itemize}

Recent domain-specific efforts---such as CyberPal.AI and SecKnowledge \cite{levi2024cyberpal,levi2025toward}---extend these principles by grounding instruction data in structured cybersecurity knowledge bases. The Athena CTI instruction framework further advances this trajectory by combining symbolic, graph-based instruction templates with benchmark-driven iteration via AthenaBench. In this sense, modern CTI-focused instruction fine-tuning can be viewed as a third-generation evolution of the original FLAN paradigm: instruction learning grounded not only in natural language, but in explicit domain structure and formal knowledge representations.


\section{Instruction Fine-Tuning and Small Language Models}
\label{sec:ift-slm}

\subsection{From Pretraining to Instruction-Following}

Pretrained LLMs are optimized for next-token prediction on broad text corpora, not for following explicit user instructions. Instruction tuning---also called instruction fine-tuning (IFT)---fine-tunes such models on labeled datasets of instructions and target outputs, improving their ability to interpret prompts as tasks rather than mere text prefixes \cite{wei2022flan,ouyang2022instructgpt}.

Key characteristics of instruction tuning include:

\begin{itemize}
  \item \textbf{Data format.} Each sample consists of an instruction, optional context, and a desired output.
  \item \textbf{Supervised objective.} The model is trained to maximize log-likelihood of the reference output given the instruction and context, often combining a standard language modeling loss with an instruction-following loss.
  \item \textbf{Generalization across tasks.} Instruction datasets deliberately mix many task types (QA, classification, summarization, transformation) so that the fine-tuned model can generalize to novel instructions, as shown by InstructGPT and FLAN \cite{wei2022flan,ouyang2022instructgpt}.
\end{itemize}

IFT is now routine in the model lifecycle: base models (e.g., LLaMA-like, Granite, Qwen) are first pretrained, then instruction-tuned, and often further refined via reinforcement learning from human feedback (RLHF) and safety fine-tunes \cite{ouyang2022instructgpt}.

\subsection{Instruction Fine-Tuning for Cybersecurity}

Cybersecurity introduces domain-specific challenges:

\begin{itemize}
  \item Technical jargon and heterogeneous schemas (CVE, CWE, CAPEC, ATT\&CK).
  \item Need for structured, evidence-grounded outputs (e.g., CVSS vectors, CWE IDs, ATT\&CK mappings).
  \item Rapid concept drift as new vulnerabilities and campaigns appear.
\end{itemize}

Several efforts address this via cyber-specific instruction data:

\begin{itemize}
  \item \textbf{SEvenLLM.} Ji et al.\ construct a bilingual cybersecurity instruction dataset and a CTI benchmark (SEvenLLM-Bench) from real incident reports, converting raw text into supervised QA, extraction, and reasoning tasks across many task types \cite{ji2024sevenllm}.
  \item \textbf{CyberBench / CyberInstruct.} Liu et al.\ create CyberInstruct, a set of cyber-specific instruction tasks drawn from logs, alerts, and vulnerability descriptions, used to train models evaluated on CyberBench \cite{liu2024cyberbench}.
  \item \textbf{Lily-Cybersecurity-7B.} Lily is a Mistral-7B fine-tune on hand-crafted cybersecurity and hacking instructions, later enriched by an LLM for style and consistency \cite{liu2024cyberbench}.
\end{itemize}

The \textbf{CyberPal.AI} line of work from IBM pushes this further by tightly integrating CTI knowledge bases and detection artifacts into an expert-driven instruction pipeline:

\begin{itemize}
  \item \textbf{SecKnowledge (CyberPal 1.0).} Levi et al.\ build an instruction dataset by mining MITRE ATT\&CK, Sigma rules, SIEM rules, adversary emulation plans, and threat reports, then using domain experts to define instruction schemas (e.g., mapping logs to ATT\&CK techniques, suggesting mitigations) and populate them semi-automatically \cite{levi2024cyberpal}.
  \item \textbf{IFT results.} CyberPal models instruction-tuned on SecKnowledge (based on LLaMA-like and other open models) outperform base models by up to roughly 24\% on security tasks and show improved robustness on CTI benchmarks \cite{levi2024cyberpal}.
\end{itemize}

The follow-on \textbf{``Toward Cybersecurity-Expert Small Language Models''} paper extends this to SecKnowledge~2.0, enriching instructions with chain-of-thought (CoT) rationales and designing a pipeline with expert-in-the-loop validation and LLM-guided knowledge retrieval, then training cybersecurity-expert SLMs (4B--20B parameter models) that rival or surpass much larger general models on CTI tasks \cite{levi2025toward}.

These works demonstrate that expert-designed, knowledge-grounded instruction datasets can transform general LLMs into specialized CTI assistants with strong reasoning and reduced hallucination---precisely the niche targeted by the Athena template-based IFT system.

\subsection{Small Language Models for Cybersecurity}

SLMs---roughly 2--20B parameters---are increasingly attractive for cybersecurity:

\begin{itemize}
  \item \textbf{Deployment constraints.} Many SOC and CTI environments require on-premises, air-gapped, or ``bring-your-own-model'' setups for regulatory and confidentiality reasons.
  \item \textbf{Latency and cost.} Incident response and detection engineering often require low-latency responses and large volumes of queries; SLMs are simpler to host and scale.
  \item \textbf{Specialization.} With strong domain IFT, SLMs can match or exceed larger general models on specialized tasks while being easier to retrain as CTI sources evolve \cite{levi2025toward}.
\end{itemize}

CyberPal~2.0 explicitly targets this: it compares cybersecurity-expert SLMs (4B--20B) against both larger open models and proprietary systems across CTI knowledge tasks, finding that well-instruction-tuned SLMs can achieve competitive or superior performance on several domains, especially when trained with CoT-enriched instructions and CTI-aware formatting \cite{levi2025toward}.

The emerging picture is that \emph{IFT + SLM + domain knowledge} is a pragmatic recipe for CTI: the heavy lifting (linguistic knowledge) comes from pretraining; domain semantics and task formats come from IFT; runtime and deployment practicality come from SLM size. The Athena IFT dataset generation system is designed to feed this recipe by generating large quantities of high-fidelity instruction triples directly from CTI knowledge bases \cite{cardei2025ift}.

\subsection{Chained SFT: Multi-Stage Supervised Fine-Tuning in Glaukopis}
\label{sec:chained-sft}

A central design choice in Glaukopis is the use of \emph{chained}, or multi-stage, supervised fine-tuning (SFT) rather than a single monolithic instruction-tuning pass. In this paradigm, the base pretrained model is adapted to its target deployment through a sequence of SFT stages, each with its own dataset, optimization hyperparameters, and intended capability outcome. The output checkpoint of each stage becomes the initialization for the next, so that capabilities accumulate compositionally rather than competing for capacity in a single fine-tune.

Chained SFT is not novel in itself: the staged structure of modern LLM post-training is well established. InstructGPT formalized a three-phase pipeline (supervised fine-tuning, reward modeling, and reinforcement learning from human feedback) in which each phase has distinct data and objectives \cite{ouyang2022instructgpt}, and FLAN/FLAN-T5 demonstrated that instruction tuning itself is most effective when it follows a well-pretrained base model rather than replacing pretraining \cite{wei2022flan,chung2022flant5}. Domain-adaptive pretraining work showed earlier that continued training of a general model on in-domain text yields measurable downstream gains, and that the order and composition of training stages materially affect final task performance \cite{gururangan2020dontstop}. In the cybersecurity setting specifically, the CyberPal/SecKnowledge~2.0 pipeline separates expert instruction-schema construction, knowledge-grounded answer synthesis, and chain-of-thought (CoT) rationale enrichment into distinct dataset-construction and training stages, with measurable benefits from the staged structure \cite{levi2024cyberpal,levi2025toward}.

Glaukopis applies this principle explicitly with three SFT stages, each producing an intermediate model that is then used as the base for the next stage. The stages are summarized below.

\paragraph{Stage 1: Foundation and domain adaptation.}
The first stage adapts a general open-weight SLM to cybersecurity language and CTI schemas. Training data consists of broad, knowledge-dense, low-structure cyber and CTI material---definitional content from MITRE ATT\&CK, CAPEC, CWE, and NVD; vendor advisories; and general-purpose security instruction data such as that surveyed in Section~\ref{sec:ift-slm}. The training objective is breadth of recall and fluent use of CTI vocabulary, rather than precise structured output. Hyperparameters favor a longer schedule, higher learning rate, and larger effective batch size suitable for absorbing wide but relatively shallow material, consistent with the domain-adaptive pretraining regime of Gururangan et al.\ \cite{gururangan2020dontstop}.

\paragraph{Stage 2: CTI task instruction tuning.}
The second stage shifts the model from general domain familiarity to precise CTI task behavior. Training data is drawn from the Athena template-based IFT generator described in Section~\ref{sec:cti-kg}: (instruction, question, answer) triples computed directly from the ASG CTI knowledge graph and covering the same task families later evaluated by AthenaBench \cite{cardei2025ift,alam2025athenabench}. Examples include CVE$\rightarrow$CWE root-cause mapping, ATT\&CK technique extraction, mitigation lookup, and platform-conditioned reasoning, with both positive and explicit negative examples derived from graph inequality constraints. Hyperparameters are tightened relative to Stage~1---lower learning rate, shorter schedule, and stricter output-format losses---reflecting the goal of crisp, schema-faithful answers rather than further broad assimilation.

\paragraph{Stage 3: Reasoning and chain-of-thought enrichment.}
The third stage trains the model to produce structured, multi-step reasoning over CTI evidence rather than terminal answers alone. Training data combines (a) CoT-augmented variants of the Stage~2 triples, in which graph-traversal reasoning chains (e.g., CVE$\rightarrow$CWE$\rightarrow$CAPEC$\rightarrow$ATT\&CK technique) are unrolled into explicit natural-language rationales, and (b) hard cases drawn from prior-stage error patterns. This stage mirrors the role of CoT enrichment in SecKnowledge~2.0, which reports that adding rationale-style supervision on top of an already CTI-tuned model improves reasoning-heavy benchmark performance more than equivalent additional non-CoT data \cite{wei2022cot,levi2025toward}. Hyperparameters use the lowest learning rate of the three stages and the shortest schedule, since the goal is to shape output style and reasoning structure on top of capabilities already present.

\paragraph{Why chain the stages.}
Two empirical motivations drive this structure. First, separating broad domain absorption from precise task behavior reduces interference: training a single model jointly on noisy domain text and structured triples tends to either underfit the structured task or overfit narrow templates, while staging allows each objective its own optimization regime \cite{gururangan2020dontstop,ouyang2022instructgpt}. Second, staging makes the pipeline incrementally refreshable in a way that matches the Athena IFT generator's incremental-generation design \cite{cardei2025ift}: when CTI sources update, only Stages~2 and~3 must be regenerated and re-run, leaving the more expensive Stage~1 foundation intact. The same property supports controlled ablation, since each intermediate checkpoint can be evaluated on AthenaBench and complementary suites (Section~\ref{sec:benchmarks}) to isolate which stage contributes which gains \cite{alam2025athenabench}.



\section{Cyber Threat Intelligence and Knowledge Graphs}
\label{sec:cti-kg}

\subsection{CTI Concepts and Workflows}

Cyber Threat Intelligence (CTI) is commonly defined as the collection, analysis, and dissemination of information about threats and adversaries to inform defensive decisions \cite{mcmillan2013gartner}. CTI workflows include:

\begin{itemize}
  \item Ingesting unstructured reports (vendor blogs, whitepapers, advisories).
  \item Structuring them into entities and relations: vulnerabilities (CVEs), weaknesses (CWEs), attack patterns (CAPEC, ATT\&CK), threat actors, campaigns, malware, tools, TTPs.
  \item Reasoning over this structure to answer questions such as:
    \begin{itemize}
      \item Which techniques does this malware implement?
      \item Which mitigations cover these techniques on a given platform?
      \item Which threat actor is most likely behind this cluster of activity?
      \item How severe and exploitable is a new CVE in our environment?
    \end{itemize}
\end{itemize}

CTI sharing and CTI knowledge graphs emphasize the importance of standardized schemas (e.g., STIX/TAXII), community knowledge bases (MITRE ATT\&CK, CAPEC, CWE), and APIs like NVD for building structured representations that can be consumed by automation \cite{strom2018mitreattck,mitreCVE,mitreCWE,mitreCAPEC}.

\subsection{MITRE Ecosystem and Related Schemas}

The MITRE ecosystem underpins much of modern CTI:

\begin{itemize}
  \item \textbf{MITRE ATT\&CK.} A globally accessible knowledge base of adversary tactics, techniques, and sub-techniques, originally released in 2013 and formally documented in ``MITRE ATT\&CK: Design and Philosophy'' \cite{strom2018mitreattck}. ATT\&CK represents adversary behavior as a graph where attack-pattern nodes (techniques) relate to software, groups (threat actors), data sources, and mitigations.
  \item \textbf{CVE and NVD.} The Common Vulnerabilities and Exposures system and NIST's National Vulnerability Database provide machine-readable vulnerability records with descriptions, affected products, CVSS metrics, and links to related CWE and CAPEC entries \cite{mitreCVE}.
  \item \textbf{CWE.} Common Weakness Enumeration captures abstract classes of software/hardware weaknesses, often linked from CVEs and CAPEC attack patterns \cite{mitreCWE}.
  \item \textbf{CAPEC.} A catalog of attack patterns describing how adversaries exploit weaknesses and how those attacks can be mitigated, with explicit links to CWE and often implicitly to ATT\&CK techniques \cite{mitreCAPEC}.
  \item \textbf{MITRE ENGAGE.} A framework for adversary engagement strategies, which can be integrated into CTI graphs as a defensive complement to ATT\&CK \cite{cardei2025ift}.
\end{itemize}

The IFT design document explicitly models these sources and their relationships in the ASG Neo4j CTI graph, extracting node types, properties, and cross-schema edges (CVE$\rightarrow$CWE, CVE$\rightarrow$CAPEC, CAPEC$\rightarrow$CWE, ATT\&CK techniques$\rightarrow$CAPEC, etc.) \cite{cardei2025ift}.

\subsection{Athena's CTI Graph and Template Language}

The \emph{Instruction Fine Tuning Dataset Design} document by Cardei defines an IFT Dataset Generation System built around an ASG CTI graph database and a rich template language \cite{cardei2025ift}. Key features include:

\begin{itemize}
  \item \textbf{CTI source support.} The system supports three categories of CTI documents: (1) structural/graph sources (MITRE ATT\&CK, CVE, CAPEC, CWE, MITRE ENGAGE, EPSS, CISA KEV); (2) LLM-assisted sources (SIGMA rules, MITRE CAR, Wazuh rules); and (3) proprietary sources (Flashpoint Ignite).
  \item \textbf{IFT triple format.} Each generated data point is an (instruction, question, answer) triple: the instruction specifies the role and style (e.g., ``You are a CTI expert that gives precise and concise answers''), the question is the user-facing CTI query, and the answer is the ground-truth response derived from the CTI graph.
  \item \textbf{Template language.} Templates provide a high-level, declarative way to bind CTI graph nodes and produce triples. They support aliases for CTI types; attribute access through nested properties; filters (using equality and inequality) applied to list properties with any/all semantics; local variables bound to nodes; random selection from lists; list-comprehension-like constructs to expand relationships into lists in the output; invisible sections that define bindings and constraints without appearing in text; and source disambiguation via a \verb|source#type| prefix (e.g., \verb|attck#attack-pattern| versus \verb|capec#attack-pattern|) \cite{cardei2025ift}.
  \item \textbf{Positive and negative examples.} The template language explicitly supports negative constraints (e.g., mitigations that do \emph{not} mitigate this technique or data components that do \emph{not} detect a technique), enabling the generation of hard negatives that teach models what does \emph{not} follow from the CTI graph.
  \item \textbf{Cross-source reasoning templates.} Example templates mimic complex reasoning chains, such as: given a CVE description, identify the corresponding CWE weakness, then find a CAPEC attack pattern with high likelihood of attack using that weakness, then find an ATT\&CK technique implementing that attack pattern.
  \item \textbf{Incremental generation.} The generator allows version and date filters at both document and concept level. Users can specify version ranges and last-modified intervals per CTI source; relationships are only used if at least one endpoint and the relationship fall within those windows.
  \item \textbf{Generation control and reporting.} A JSON job configuration file specifies CTI version ranges, template texts, per-template count limits and coverage fractions, output directories, and dataset versioning. The generation API produces both per-template triple files and a global report summarizing counts, coverage, and errors, which can be post-processed into human-readable formats \cite{cardei2025ift}.
\end{itemize}

Conceptually, this system is a symbolic counterpart to the LLM-assisted instruction pipelines used in CyberPal and SEvenLLM: where CyberPal uses LLMs and experts to synthesize security instructions from unstructured sources, the Athena system uses a declarative template language over a CTI graph to generate strongly grounded and precisely controllable instruction triples.

\subsection{CTI Knowledge Graphs and LLMs}

Recent work shows that LLMs and CTI knowledge graphs are synergistic. Various efforts use LLMs to extract triples from CTI text and populate knowledge graphs, then perform link prediction and reasoning over those graphs to produce actionable recommendations \cite{fieblinger2024kg,zhang2023cyberkg,wu2022opencykg,yu2024ctikg,zhao2024kctiaa}. Other systems build CTI-KGs directly from ATT\&CK, CVE, threat reports, and open sources, then query them using LLMs or hybrid systems for tasks such as threat hunting and reporting.

The Athena IFT design follows a complementary pattern: \emph{knowledge graph $\rightarrow$ templates $\rightarrow$ IFT triples $\rightarrow$ SLM}, with AthenaBench serving as a \emph{knowledge graph $\rightarrow$ evaluation tasks} pipeline. This alignment is central to the next section.

\section{Cybersecurity LLM Benchmarks and AthenaBench}
\label{sec:benchmarks}

\subsection{Why Cybersecurity-Specific Benchmarks?}

General benchmarks such as GLUE, MMLU, and BIG-Bench measure broad language understanding and world knowledge, and certain cybersecurity-adjacent capabilities (e.g., information-security multiple-choice questions in MMLU) appear as minor subsets, but these suites are not designed to probe the workflows that actually define security practice. They lack the structured schemas, evidence-grounded reasoning, and threat-current content that distinguish cybersecurity tasks from general knowledge work. The gaps that motivate domain-specific benchmarking are concrete:

\begin{itemize}
  \item \textbf{Schema fidelity.} Security workflows operate over standardized identifiers and structured artifacts (CVE, CWE, CAPEC, CVSS vectors, MITRE ATT\&CK technique IDs, STIX/TAXII objects). General benchmarks do not test whether a model produces well-formed outputs in these schemas, which is the difference between a usable assistant and one that requires constant human re-formatting.
  \item \textbf{Applied reasoning.} CTI and SOC analysts rarely need pure recall; they need chained inferences (e.g., from a CVE description to the underlying weakness class, then to relevant attack patterns and detection opportunities), severity estimation under partial information, and abductive attribution from observed TTPs. These reasoning shapes are largely absent from general benchmarks.
  \item \textbf{Temporal relevance.} The threat landscape evolves on a weekly cadence (new CVEs, new APT activity, new exploitation campaigns), so static benchmarks rapidly drift from operational reality and become measures of stale memorization rather than current competence \cite{alam2025athenabench}.
  \item \textbf{Data contamination.} Many cybersecurity sources (MITRE ATT\&CK, the NVD, public threat reports) are heavily represented in LLM pretraining corpora. A static benchmark drawn from those same sources risks measuring recall of training data rather than genuine reasoning, an effect that dynamic benchmarks aim to mitigate \cite{alam2025athenabench}.
\end{itemize}

In response, the past two years have seen a rapid expansion of cybersecurity-specific benchmarks, which can be grouped roughly into four overlapping families: \emph{knowledge benchmarks}, focused on factual recall of security concepts and terminology; \emph{secure-coding and offensive-capability benchmarks}, focused on whether models help or hinder secure software development and adversarial use; \emph{CTI and SOC benchmarks}, focused on applied threat-intelligence and operations reasoning; and \emph{dynamic / live-source benchmarks}, focused on resistance to data staleness and contamination. Important entries include:

\begin{itemize}
  \item \textbf{CyberMetric.} A retrieval-augmented, certification-style multiple-choice benchmark by Tihanyi et al.\ covering nine cybersecurity domains, released in graduated sizes (80, 500, 2000, and 10\,000 questions) to support both rapid evaluation and statistically robust comparison \cite{tihanyi2024cybermetric}.
  \item \textbf{SecEval.} A multiple-choice benchmark of cybersecurity knowledge questions spanning multiple security subdomains, used to evaluate factual coverage of foundation models on standard certification-style content \cite{hu2024seceval}.
  \item \textbf{CyberBench.} A multi-task benchmark for cybersecurity LLMs that includes classification, question answering, and log/report analysis, accompanied by the CyberInstruct training corpus used by the same authors to build cyber-tuned models \cite{liu2024cyberbench}.
  \item \textbf{CTIBench.} A dedicated CTI benchmark covering CTI knowledge, root-cause mapping, vulnerability severity prediction, threat-actor attribution, threat-knowledge graph construction, and attack-technique extraction; constructed from curated CTI corpora and designed for reasoning-intensive evaluation \cite{alam2024ctibench}.
  \item \textbf{SEvenLLM-Bench.} A bilingual CTI benchmark derived from real incident reports, designed alongside the SEvenLLM-Instruct dataset and emphasizing cross-lingual security comprehension \cite{ji2024sevenllm}.
  \item \textbf{Purple Llama CyberSecEval and CyberSecEval~2.} Benchmarks from Meta evaluating secure-code generation, susceptibility to assisting with cyberattacks, prompt-injection robustness, and related risk dimensions \cite{bhatt2023purplellama,bhatt2024cyberseceval2}.
  \item \textbf{WMDP.} The Weapons of Mass Destruction Proxy benchmark by Li et al.\ includes a cybersecurity subset measuring potentially hazardous cyber-offense knowledge and is widely used to study unlearning and safety-capability tradeoffs in security-adjacent capabilities \cite{li2024wmdp}.
  \item \textbf{AthenaBench.} A dynamic CTI benchmark that draws from live sources (MITRE ATT\&CK, NVD, CISA advisories, threat reports) and explicitly targets the staleness and contamination problems that limit static suites; discussed in detail below \cite{alam2025athenabench}.
\end{itemize}

Together, these benchmarks shift evaluation from generic question answering toward security-aware reasoning, but they cover very different surfaces---factual recall, secure-code generation, applied CTI reasoning, and dynamic resistance to drift---and no single benchmark is sufficient on its own. Glaukopis is therefore evaluated against a portfolio of complementary suites, with particular weight on CTIBench and AthenaBench (which directly target CTI reasoning), CyberMetric at the 2\,000- and 10\,000-question tiers (which provides statistically robust factual coverage), and CyberSOCEval's malware and threat-intelligence test suites (which probe operationally relevant SOC reasoning). The remainder of this section describes these benchmarks in more detail.

\subsection{CyberMetric}
\label{sec:cybermetric}

CyberMetric is a multiple-choice cybersecurity knowledge benchmark constructed by Tihanyi et al.\ using a retrieval-augmented generation (RAG) pipeline over an authoritative corpus of cybersecurity standards, frameworks, NIST publications, and academic and practitioner references \cite{tihanyi2024cybermetric}. Candidate questions are generated by an LLM conditioned on retrieved source material, then filtered and validated by human cybersecurity experts before inclusion in the released dataset. The result is a benchmark whose questions are grounded in citable source material rather than constructed from scratch, which mitigates the common problem of LLM-as-author benchmarks whose answer keys reflect the generating model's biases more than the underlying domain.

CyberMetric covers nine cybersecurity domains: network security, cryptography, identity and access management, security policies and governance, threats and vulnerabilities, security operations, application and system security, cloud security, and incident response. It is released in four cumulative sizes---80, 500, 2\,000, and 10\,000 questions---designed to support different evaluation regimes: the 80- and 500-question tiers for rapid iteration and ablation, and the 2\,000- and 10\,000-question tiers for statistically robust comparison between models where small accuracy differences must be distinguished from sampling noise \cite{tihanyi2024cybermetric}.

In Glaukopis, CyberMetric is used at the 2\,000- and 10\,000-question tiers. These tiers were chosen for three reasons. First, sample size: at 10\,000 questions the standard error on accuracy estimates is approximately 0.5 percentage points, which is fine-grained enough to detect the kind of single-digit improvements typically produced by an additional SFT stage (Section~\ref{sec:chained-sft}). Second, domain coverage: the larger tiers include enough questions per domain to support per-domain breakdowns, which lets us identify whether a given training-stage change improves balanced cybersecurity knowledge or only narrow subdomains. Third, comparability: published evaluations of contemporary models---both general-purpose frontier models and cyber-specialized SLMs---report CyberMetric scores at the 2\,000- and 10\,000-question tiers, allowing direct comparison without re-running baselines.

CyberMetric's principal limitation, shared with all multiple-choice benchmarks, is that it measures recognition rather than generation: a model can score well on CyberMetric without being able to produce structured CTI outputs or coherent multi-step reasoning. For this reason it is used in Glaukopis as one component of a broader evaluation portfolio rather than as a primary metric, and is paired with reasoning-oriented suites (CTIBench, AthenaBench) and operationally focused suites (CyberSOCEval) that exercise generation and applied reasoning.

\subsection{CyberSOCEval}
\label{sec:cybersoceval}

%% TODO(author): replace cybersoceval2025 with the exact citation key
%% used in your converted glaukopis.bib for the CyberSOCEval reference.
CyberSOCEval is an evaluation suite oriented around security-operations-center (SOC) tasks, providing test suites that target operationally meaningful reasoning rather than pure knowledge recall \cite{cybersoceval2025}. Whereas CyberMetric measures what a model knows about cybersecurity as a body of knowledge, and CTIBench/AthenaBench measure what a model can reason over a CTI knowledge graph, CyberSOCEval measures whether a model can perform the kind of analytic work that occupies an analyst inside a SOC: classifying observed artifacts, correlating evidence, and producing actionable summaries.

In Glaukopis, CyberSOCEval is used along two of its test suites:

\begin{itemize}
  \item \textbf{Malware analysis suite.} Tasks in this suite present malware-related artifacts (samples, behavioral indicators, or analyst-style descriptions) and probe whether the model can extract family/variant attribution, characterize behavior in a structured way (e.g., what techniques the artifact exhibits), and identify defensive implications. This suite is operationally important because malware-related reasoning is one of the most frequent generation tasks asked of CTI assistants and is also one of the most failure-prone: errors here typically take the form of confident misattribution rather than refusal, which is more harmful than abstention.
  \item \textbf{Threat-intelligence suite.} Tasks here exercise the analyst-facing CTI reasoning chain: ingesting raw report or alert content, extracting TTPs and indicators, mapping them to standardized schemas (ATT\&CK techniques, mitigations, threat-actor designations where supported by evidence), and producing structured CTI summaries. This complements AthenaBench's task structure but emphasizes free-form narrative input over knowledge-graph-derived queries, which makes it a useful test of how well a CTI SLM generalizes from training data structured by the Athena template generator to less-structured real-world inputs.
\end{itemize}

These two suites are particularly relevant to Glaukopis because they exercise capabilities that the chained-SFT pipeline is explicitly designed to develop. Stage~2 (CTI task SFT, Section~\ref{sec:chained-sft}) trains the model on schema-faithful answers to graph-derived questions, and Stage~3 (CoT enrichment) trains it on multi-step reasoning traces; the malware and threat-intelligence suites of CyberSOCEval test whether those capabilities transfer to operationally realistic narrative inputs. Their inclusion alongside CyberMetric (factual coverage), CTIBench (CTI reasoning), and AthenaBench (dynamic CTI reasoning) gives the evaluation portfolio coverage across the four benchmark families identified above: knowledge, secure-coding/offensive capability via Purple Llama suites as needed, applied CTI reasoning, and operational SOC reasoning \cite{tihanyi2024cybermetric,alam2024ctibench,alam2025athenabench,bhatt2023purplellama,bhatt2024cyberseceval2,cybersoceval2025}.

\subsection{AthenaBench: A Dynamic CTI Benchmark}

AthenaBench is introduced as a dynamic benchmark suite for evaluating LLMs on realistic, knowledge-intensive CTI tasks and explicitly extends CTIBench \cite{alam2025athenabench,alam2024ctibench}. Its key innovations include:

\begin{itemize}
  \item \textbf{Dynamic data construction.} Instead of static corpora, AthenaBench connects to live CTI data sources and APIs (MITRE ATT\&CK, NVD, CISA advisories, threat reports) to continuously generate benchmark samples within configurable time windows. Tasks such as root cause mapping (CVE$\rightarrow$CWE) and severity prediction (CVSS) are drawn from recent CVEs and weaknesses, mitigating training-set leakage and static knowledge effects \cite{alam2025athenabench}.
  \item \textbf{Six complementary tasks.} AthenaBench defines: CTI Knowledge Test (CKT), Root Cause Mapping (RCM), Vulnerability Severity Prediction (VSP), Threat Actor Attribution (TAA), Risk Mitigation Strategy (RMS), and Attack Technique Extraction (ATE). Together, they probe CTI knowledge, root cause reasoning, structured severity prediction, abductive attribution, defensive planning, and TTP extraction \cite{alam2025athenabench}.
  \item \textbf{Enhanced metrics and aggregated score.} AthenaBench improves task-specific metrics and defines a unified aggregated score, making it easier to compare LLMs' CTI reasoning capabilities along different dimensions \cite{alam2025athenabench}.
  \item \textbf{Model evaluation.} The authors evaluate 12 LLMs, including GPT-4, GPT-4o, GPT-5, Gemini-2.5 Flash/Pro, Qwen variants, Llama 3/3.1/3.3, and a Llama-Primus merged model. Proprietary models (GPT-5, Gemini-2.5 Pro) lead overall, but all models struggle on reasoning-intensive tasks such as TAA and RMS, especially open-source ones \cite{alam2025athenabench}.
\end{itemize}

By using live CTI sources and emphasizing reasoning over current threats, AthenaBench moves CTI benchmarking closer to real-world analyst workloads and provides a stringent target for cybersecurity-expert SLMs.

\subsection{Relationship to Athena's IFT Dataset Generator}

AthenaBench and the Athena IFT generator are complementary components of a full CTI LLM lifecycle:

\begin{itemize}
  \item \textbf{Training side (IFT generator).} Uses CTI graph data (ATT\&CK, CVE, CAPEC, CWE, ENGAGE, etc.) and a declarative template language to generate instruction triples that drill CTI knowledge, multi-hop reasoning, negative cases, and structured outputs \cite{cardei2025ift}.
  \item \textbf{Evaluation side (AthenaBench).} Dynamically constructs CTI tasks from the same or similar underlying sources (ATT\&CK, NVD, CISA, threat reports) but with different pipelines and sampling strategies, providing independent, evolving test sets to evaluate model generalization \cite{alam2025athenabench}.
\end{itemize}

Many of the template patterns in the design document directly mirror AthenaBench tasks:
\begin{itemize}
  \item CVE$\rightarrow$CWE reasoning templates align with RCM.
  \item Templates that traverse CVE$\rightarrow$CWE$\rightarrow$CAPEC$\rightarrow$ATT\&CK map onto knowledge and mitigation reasoning tested in CKT and RMS.
  \item Templates involving malware, techniques, platforms, and data components parallel the kind of ``who uses what and how is it detected'' reasoning central to CTI analysis and attack technique extraction tasks \cite{cardei2025ift}.
\end{itemize}

The incremental generation features of the IFT system (version and timestamp filtering) make it possible to align the training data's recency with AthenaBench's dynamic evaluation horizon, mitigating data leakage while ensuring models are trained on approximately the same vintage of CTI as they are evaluated on \cite{cardei2025ift,alam2025athenabench}.

This alignment positions the Athena IFT generator as a natural training counterpart to AthenaBench's evaluation pipeline, much as SecKnowledge and CyberPal models are paired with CTI benchmarks in IBM's work \cite{levi2024cyberpal,levi2025toward}.

\section{Positioning Template-Based IFT for Cybersecurity SLMs}
\label{sec:athena-ift}

\subsection{Design Goals and Advantages}

Compared to purely LLM-generated instruction datasets, a template-driven, knowledge-grounded approach offers several advantages:

\begin{itemize}
  \item \textbf{Ground-truth faithfulness.} Answers are computed directly from an authoritative CTI graph, reducing hallucinations and inconsistencies in training labels.
  \item \textbf{Explicit coverage control.} Templates can be designed to cover specific combinations of entities (e.g., CAPEC patterns affecting memory resources and their detection methods) or relationships (e.g., malware using two techniques on Linux), and per-template coverage fractions and count limits control dataset size and diversity \cite{cardei2025ift}.
  \item \textbf{Rich negative examples.} Inequality constraints and relationship filters enable systematic generation of negative examples such as mitigations that do not address a technique or data components that do not detect any technique used by a malware family.
  \item \textbf{Explainability of templates.} Each template corresponds to a human-readable CTI reasoning pattern, making it easier for security experts to audit and refine the dataset.
  \item \textbf{Incremental refresh.} Version and date filters allow periodic regeneration of updated IFT corpora as CTI sources evolve, paralleling AthenaBench's dynamic evaluation but on the training side \cite{cardei2025ift,alam2025athenabench}.
\end{itemize}

In combination with SLMs, this supports a training loop where:

\begin{itemize}
  \item New CTI sources or schema updates lead to a refreshed CTI graph.
  \item CTI experts design or revise templates for new patterns and tasks.
  \item The generator produces updated IFT triples for fine-tuning SLMs.
  \item Models are evaluated on AthenaBench and other benchmarks.
  \item Discrepancies and weaknesses inform template revisions and possibly benchmark extensions.
\end{itemize}

\subsection{Comparison with IBM's CyberPal and Other Instruction Pipelines}

The Athena approach is conceptually closely related to SecKnowledge/CyberPal, but with different emphasis:

\begin{itemize}
  \item \textbf{CyberPal / SecKnowledge.} Use expert-defined schemas and a pipeline combining LLMs and scripts to synthesize instructions from CTI sources and detection artifacts (e.g., Sigma rules). Emphasize human-curated ``expert instruction patterns'' and LLM-assisted generation, including CoT answers in SecKnowledge~2.0 \cite{levi2024cyberpal,levi2025toward}.
  \item \textbf{Athena IFT generator.} Uses a symbolic template language compiled into CTI graph queries, producing triples that are guaranteed to be consistent with the knowledge graph, with strong use of cross-schema constraints and negative examples, and built around incremental, configuration-driven generation \cite{cardei2025ift}.
\end{itemize}

Both approaches are complementary:

\begin{itemize}
  \item CyberPal demonstrates the effectiveness of expert instruction schemas and CoT enrichment in achieving high cybersecurity performance, especially in SLMs.
  \item Athena's system provides precise and reproducible control over content generation aligned with CTI schemas and knowledge graph structure, potentially feeding into similar CoT-enriched pipelines (e.g., by later adding CoT rationales using LLMs on top of graph-derived answers).
\end{itemize}

\subsection{Toward Cybersecurity-Expert Small Language Models}

Given these ingredients, an end-to-end cybersecurity-expert SLM pipeline could look like:

\begin{enumerate}
  \item \textbf{CTI graph construction.} Continuously ingest and normalize CTI sources (ATT\&CK, CVE, CAPEC, CWE, ENGAGE, ICS/OT sources, proprietary feeds) into a Neo4j or similar graph.
  \item \textbf{Template authoring.} CTI experts encode important reasoning patterns (knowledge recall, root cause mapping, severity estimation, attribution clues, mitigation reasoning) as templates in the Athena language.
  \item \textbf{IFT generation.} Run the generation tool with version and date constraints aligned to a training window, producing large, structured IFT datasets per task family \cite{cardei2025ift}.
  \item \textbf{Model training.} Fine-tune SLMs (e.g., 4--14B open models) using these IFT datasets, optionally adding CoT rationales, conversation styles, and alignment methods similar to CyberPal SecKnowledge~2.0 \cite{levi2025toward}.
  \item \textbf{Evaluation.} Benchmark models on AthenaBench and complementary suites (CTIBench, SecEval, CyberBench, SEvenLLM-Bench) to assess generalization and compare with larger proprietary models \cite{alam2024ctibench,alam2025athenabench,hu2024seceval,liu2024cyberbench,ji2024sevenllm}.
  \item \textbf{Feedback and iteration.} Use benchmark error patterns to design new templates, adjust CTI graph coverage, or refine IFT data (e.g., adding more negative cases or multi-hop chains).
\end{enumerate}

This is closely aligned with IBM's vision of cybersecurity-expert SLMs but anchored in Athena's graph-driven IFT and dynamic CTI benchmarking.

\section{Open Challenges and Research Directions}

Despite rapid progress, several open questions remain:

\begin{itemize}
  \item \textbf{Data leakage and benchmark freshness.} Dynamic benchmarks like AthenaBench reduce, but do not eliminate, the risk that training data (e.g., from the same CTI sources) overlaps heavily with evaluation samples. Careful time-based splits and provenance tracking between the IFT generator and benchmarks will be crucial \cite{alam2025athenabench,cardei2025ift}.
  \item \textbf{Balancing symbolic templates and LLM creativity.} Template-based generation offers precision and coverage, but LLM-generated paraphrases can increase linguistic diversity and realism (e.g., multiple phrasings of the same CTI scenario). Hybrid pipelines that use templates to define semantics and LLMs to diversify surface form---while preserving correctness---are a promising area.
  \item \textbf{Incorporating chain-of-thought and explanations.} SecKnowledge~2.0 shows that CoT rationales significantly improve cyber reasoning \cite{levi2025toward}. Extending Athena templates to generate structured reasoning steps (via explicit graph traversal reasoning) and then augmenting them with LLM-generated natural language CoT could yield explainable CTI SLMs.
  \item \textbf{Operational evaluation dimensions.} AthenaBench focuses on knowledge and reasoning quality. Security operations also care about robustness to prompt injection, misuse, adversarial CTI, latency, and integration with SIEM/SOAR systems. Benchmarks like SecEval and CyberBench partly address this; extending AthenaBench and complementary suites to cover these aspects remains an important direction \cite{hu2024seceval,liu2024cyberbench,alam2025athenabench}.
  \item \textbf{Multilingual and cross-regional CTI.} SEvenLLM emphasizes bilingual corpora \cite{ji2024sevenllm}. Current CTI benchmarks primarily use English sources, omitting intelligence from other languages that often contain region-specific threats. Template-based IFT over multilingual CTI graphs and multilingual benchmark tasks could help close this gap.
  \item \textbf{Human-in-the-loop and analyst trust.} For high-stakes CTI decisions, analysts need to understand why the model suggested a mapping or mitigation and how it connects to CTI sources. Template-driven data, CTI graph citations, and CoT rationales may enable verifiable, inspectable explanations.
\end{itemize}

\section{Conclusion}

Instruction fine-tuning and small language models are emerging as a practical and powerful combination for cybersecurity, especially in CTI workflows that demand structured reasoning, up-to-date knowledge, and on-premises deployment. IBM's CyberPal and SecKnowledge series demonstrate that expert-curated, CTI-grounded instruction datasets can turn SLMs into strong cybersecurity assistants, while AthenaBench raises the bar for CTI evaluation by integrating live CTI sources and focusing on reasoning-intensive tasks \cite{levi2024cyberpal,levi2025toward,alam2025athenabench}.

The Athena IFT Dataset Generation System designed by Cardei and collaborators provides a complementary training-time infrastructure: a template language over a CTI knowledge graph that can generate rich, grounded, and incrementally refreshable instruction triples across multiple CTI schemas and tasks \cite{cardei2025ift}. Together, these elements outline a coherent research program:

\begin{itemize}
  \item Build and maintain high-fidelity CTI graphs from authoritative sources.
  \item Encode analyst reasoning patterns as templates that produce IFT data.
  \item Train cybersecurity-expert SLMs using these datasets.
  \item Evaluate them with dynamic CTI benchmarks like AthenaBench and iteratively refine both data and models.
\end{itemize}

The result is a principled pathway toward cybersecurity-expert small language models grounded simultaneously in CTI knowledge bases, instruction fine-tuning theory, and rigorous CTI benchmarking.


\section{Conflicts of interest}
None declared.

\section{Funding}
No external funding supported this manuscript.

\section{Data availability}
No datasets were generated or analyzed during the preparation of this survey paper.

\section{Acknowledgments}
The authors thank colleagues at Florida Atlantic University's Department of Electrical Engineering and Computer Science and at Athena Security Group for discussions that informed this work.


\bibliographystyle{abbrvnat}
\bibliography{glaukopis}

\end{document}
